\section{Related Work} \label{sec:related}

Security studies in autonomous drones, cars, and mobile robots often
evaluate one attack or one defense in isolation.  Prior work includes
(i)~fault/attack injection methods from dependable-systems research and
(ii)~intrusion detection and runtime monitoring approaches for CPS/IoT.
However, experiment setups are frequently ad hoc, artifacts and datasets
are not consistently shared, and evaluations are difficult to reproduce
or compare across studies.

\subsection{CPS and Autonomous System Security}

Shoukry et~al.~\cite{shoukry2013} demonstrated GPS spoofing attacks on
autonomous vehicles and proposed sensor-fusion countermeasures, but
evaluated a single attack--defense pair.  Mo and Sinopoli~\cite{mo2009}
analyzed replay attacks on cyber-physical control systems with a focus on
detection theory rather than a runnable evaluation framework.
Neshenko et~al.~\cite{neshenko2019} surveyed IoT security datasets and
metrics, identifying a lack of standardized evaluation environments.

More recently, the SENTINEL survey~\cite{sentinel2026} emphasizes that
detection alone is insufficient in safety-critical CPS and proposes a
lifecycle-aware threat methodology for AI agents, directly supporting
the need for integrated evaluation frameworks that measure both security
and mission-level outcomes.  Work on automotive IDS
evasion~\cite{autoids2025} shows that CAN-bus intrusion detectors can be
evaded across white-box, grey-box, and black-box scenarios, reinforcing
the argument that single-defense evaluation is insufficient.

\subsection{Fault Injection and Testing Frameworks}

FIMSim~\cite{fimsim} provides fault injection for multi-robot systems
using a simulator-integrated approach, but focuses on Byzantine faults
rather than adversarial security attacks and does not include defense
modules.  Rodrigues et~al.~\cite{rodrigues2020} and
Machado et~al.~\cite{machado2021} contribute fault injection
methodologies for robotic platforms, but neither provides a pluggable
defense interface or cross-comparison infrastructure.

MIXED-SENSE~\cite{mixedsense2025} is the closest domain match, offering
a mixed-reality sensor emulation framework for UAV testing that can
benchmark attack detector performance.  However, it handles only one
attack type (false data injection) and does not support arbitrary catalogs
of attacks and defenses.

\subsection{Security Evaluation Benchmarks}

CARE~\cite{care2024} provides the most comprehensive cross-evaluation to
date, comparing 15 attacks against 8 defenses across 12 network IDS
detectors.  It demonstrates that cross-defense comparison is valuable
for identifying blind spots.  However, CARE operates entirely on generic
network security datasets (malware, spam, network intrusion) and has no
simulator, no autonomous-system telemetry, and no mission context.

DYNAMITE~\cite{dynamite2025} uses ML to select the best defense per
attack type for IoT network traffic, showing a 96.2\% reduction in
compute time and 76.7\% F1 improvement over random defense selection.
This motivates multi-defense comparison but targets flat IoT traffic,
not autonomous system command-and-control streams.

SAGE~\cite{sage2025} proposes attack-aware adaptive defense selection,
achieving a 201\% F1 improvement over the state of the art and narrowing
the gap to an oracle selector to 3.8\%.  This validates the argument
that adaptive and attack-aware defenses outperform static approaches.

UAVIDS-2025~\cite{uavids2025} contributes 122,171 labeled flow records
across five attack types with 22 features per flow, demonstrating demand
for UAV IDS evaluation resources.  However, it is a static dataset only:
no runnable orchestrator, no pluggable defenses, no synchronized
real-time logging, and no mission metrics.

IForestExplain~\cite{iforestexplain2025} applies an Isolation Forest with
explainability to IoT anomaly detection, showing that correlation-based
feature selection outperforms other approaches.  This validates our
technical choice of Isolation Forest for CPS anomaly detection, though
their work targets generic IoT traffic rather than autonomous systems.

\subsection{Summary of Differences}

Our work differs from the above in that we provide a \emph{runnable},
\emph{extensible} evaluation harness specifically designed for
autonomous/CPS platforms.  Unlike CARE, we include a protocol simulator
and mission-level metrics.  Unlike MIXED-SENSE, we support an arbitrary
catalog of attacks and defenses.  Unlike UAVIDS-2025, we provide an
orchestrated execution environment rather than a static dataset.  The
comparison in Table~\ref{tab:comparison} shows that no prior framework
simultaneously covers all five dimensions: multi-attack, multi-defense,
CPS/autonomous context, runnable harness, and mission metrics.
